\begin{definition}[$\tilde\psi_{m,Q,\epsilon,C}$]
  Let $E$ be a metric space, $\gamma \in \Theta(E)$ (\noderef{Curve}), $m \in \mathbb{N}$,
  $Q \subset E$ a finite set, and $\epsilon, C \in \mathbb{R}$. Write $s_i := i \cdot
  \length(\gamma)/m$ for $i = 0,\dotsc,m$. Define
  \[
    \tilde\psi_{m,Q,\epsilon,C}(\gamma) := 2C^2\frac{\length(\gamma)^2}{m}
      + C\frac{\length(\gamma)}{m} \sum_{i=1}^m \min\left\{2C,\; \epsilon+
        2Cd(\gamma(s_i),Q)\right\}
    ,
  \]
  where $d(x,Q) = \inf_{q \in Q} d(x,q)$ (Mathlib's \texttt{Metric.infDist}, with the convention
  $d(x,\emptyset) = 0$). This is exactly the second display of paper Lemma~A.2 (paper.tex line
  1722--1723), with the same widening of the quantification over $Q,\epsilon,C$ to arbitrary finite
  $Q \subset E$ and real $\epsilon, C$ as in \noderef{psi}, for the same reason (the paper's
  restriction to rationals and to a fixed countable dense $E_{\mathbb{Q}}$ is a countability device
  for the Strong Law step, not part of Lemma~A.2's or Corollary~A.6's mathematical content).
\end{definition}
